\section{방법론}
\label{sec:methods}

\subsection{PUBG 텔레메트리의 시공간 데이터베이스 표현}

PUBG 공식 텔레메트리는 원래 시간순 이벤트 로그로 제공된다. 이 로그를
매치 종료 후의 집계 행으로 사용하지 않고, 페이즈별 생존 예측을 위한 시공간
데이터베이스 표현으로 재구성한다. 핵심은 팀을 이동 객체로 두고, 안전구역과
블루존을 시간에 따라 변하는 영역으로 두며, 피해 사건, 기절 사건, 처치 사건, 탈락
사건을 시공간 사건으로 정렬하는 것이다. 그리고 각 페이즈 종료 시점까지 살아남은
팀만 다음 페이즈의 예측 단위로 남긴다.

\begin{table*}[t]
\centering
\caption{PUBG 텔레메트리의 시공간 데이터베이스 구성요소}
\label{tab:stdb-components}
\small
\setlength{\tabcolsep}{3pt}
\begin{tabular}{p{0.17\linewidth}p{0.22\linewidth}p{0.35\linewidth}p{0.17\linewidth}}
\toprule
시공간 DB 구성요소 & PUBG 텔레메트리에서의 대응 & 생성 방식 & 생존 예측에서의 역할 \\
\midrule
이동 객체
& 스쿼드 팀의 중심 위치
& 10초 간격 팀 중심 좌표 \(x,y,z\)와 속도 계산
& 이동 이력 \\
시간변화 영역
& 안전구역, 블루존
& 페이즈별 중심, 반경, 경계 거리, 내부 여부 추출
& 공간 위험 맥락 \\
시간변화 속성
& 체력, 생존 인원, 차량 이용 여부
& 페이즈 절단 시점 이전의 팀 상태 집계
& 팀 생존 여력 \\
시공간 사건
& 피해 사건, 기절 사건, 처치 사건, 탈락 사건
& 시간창 내 사건을 발생 시점과 관련 팀 단위로 정렬
& 교전 압력 및 라벨 생성 \\
위험집합
& 페이즈 \(k\)까지 생존한 팀
& 다음 페이즈 탈락 여부를 팀-페이즈 단위로 라벨링
& 예측 단위 \\
\bottomrule
\end{tabular}
\end{table*}

표~\ref{tab:stdb-components}의 목적은 단순한 변수 목록을 제시하는 데 있지 않다.
이 표는 PUBG 텔레메트리가 최종 순위 예측을 위한 정적 입력이 아니라, 이동 객체와
시간변화 영역, 사건, 위험집합이 결합된 시공간 데이터베이스로 재구성될 수
있음을 정리한다. 이후의 예비 실험은 이 구성요소들이 페이즈별 생존 예측에서 어떤
방식으로 작동하는지 확인하기 위한 것이다.

\subsection{페이즈별 위험집합과 예측 라벨 구성}

본 연구의 분석 단위는 매치 전체가 아니라 페이즈별 생존 팀이다. 팀 \(i\)가
페이즈 \(k\)의 종료 시점까지 생존한 경우, 해당 팀은 다음 페이즈의 위험집합에
포함된다. 이때 목표값은 다음 페이즈에서 탈락했는지 여부로 정의한다.
\[
y_{i,k} =
\begin{cases}
1, & \text{팀 } i \text{가 페이즈 } k+1 \text{에서 탈락한 경우},\\
0, & \text{그렇지 않은 경우}.
\end{cases}
\]
입력 \(X_{i,k}\)는 페이즈 \(k\)의 종료 시점까지 관측된 텔레메트리만 포함한다.
예측 시점 이후의 피해량, 위치, 탈락 사건은 입력에서 제외하였다. 이는 모델이 사후
정보를 이용해 결과를 확인하는 것을 방지하기 위한 설정이다.

이 정의에서 예측 시점은 항상 한 페이즈 앞선다. 예를 들어 한 팀이 페이즈 1의
종료 시점까지 생존했다면, 모델은 페이즈 1 종료 이전의 10초 단위 스냅샷 이력만
사용하여 페이즈 2에서 탈락할지를 예측한다. 페이즈 2 안에서 발생한 이동, 피해,
탈락 사건은 이 예측에 들어가지 않는다. 마찬가지로 페이즈 2 종료 시점까지 생존한
팀은 그 시점 이전의 이력으로 페이즈 3의 탈락 여부를 예측한다.

\subsection{시공간 표현 요소의 추출}

아래 두 표가 본 실험의 feature 정의다. Raw telemetry를 직접 모델에 넣은 것이 아니라,
먼저 10초 간격의 팀-스냅샷 node tensor \(X_s\)로 변환한 뒤, 조건별로 feature block을
나누거나 다시 조합하여 학습했다.

\begin{table*}[t]
\centering
\caption{Raw telemetry에서 구성한 모델 입력 feature block}
\label{tab:raw-feature-blocks}
\scriptsize
\setlength{\tabcolsep}{3pt}
\begin{tabular}{L{0.18\linewidth}L{0.42\linewidth}L{0.24\linewidth}L{0.10\linewidth}}
\toprule
입력 block & 모델 입력 column & 생성 방식 & 실험 역할 \\
\midrule
이동
& \texttt{pos\_x}, \texttt{pos\_y}, \texttt{pos\_z}, \texttt{speed\_mps}
& 팀 중심 좌표와 직전 스냅샷 대비 속도
& 이동 이력 \\
안전구역
& \texttt{rel\_zone\_x}, \texttt{rel\_zone\_y}, \texttt{dist\_blue\_boundary\_m}, \texttt{inside\_blue}, \texttt{dist\_target\_center\_m}
& safe zone/blue zone에 대한 상대 위치와 경계 거리
& zone 맥락 \\
생존 상태
& \texttt{mean\_health}, \texttt{min\_health}, \texttt{alive\_n}
& 플레이어 상태를 팀 단위로 집계
& 생존 여력 \\
자원 상태
& \texttt{vehicle\_ratio}
& 팀원 중 차량 탑승 비율
& 이동 자원 \\
\bottomrule
\end{tabular}
\end{table*}

\begin{table*}[t]
\centering
\caption{Feature block을 분해·조합한 학습 조건}
\label{tab:feature-condition-map}
\small
\setlength{\tabcolsep}{4pt}
\begin{tabular}{L{0.28\linewidth}L{0.42\linewidth}L{0.22\linewidth}}
\toprule
학습 조건 & 사용한 \(X_s\) 조합 & 시간 처리 \\
\midrule
최신 상태 기준
& 13개 전체, 단 최신 스냅샷만 사용
& 각 phase cut의 마지막 snapshot만 사용 \\
팀 생존 여력
& 생존 상태 + 자원 상태
& phase cut 이전 10초 스냅샷 시퀀스 \\
이동 이력
& 이동 block
& phase cut 이전 10초 스냅샷 시퀀스 \\
안전구역 맥락
& 안전구역 block
& phase cut 이전 10초 스냅샷 시퀀스 \\
이동+안전구역
& 이동 block + 안전구역 block
& phase cut 이전 10초 스냅샷 시퀀스 \\
전체 스냅샷
& 13개 전체의 10초 이력
& phase cut 이전 10초 스냅샷 시퀀스 \\
\bottomrule
\end{tabular}
\end{table*}

즉 본 실험은 \(X_s\)를 block 단위로 나눈 ablation이다. 최신 상태 기준은 각 phase
cut에서 마지막으로 관측된 snapshot만 사용하고, 나머지 시간 이력 조건은 10초 간격
snapshot sequence를 평균이나 합계로 먼저 뭉개지 않는다. 모델은 이 순서를 인과적
시간 합성곱으로 인코딩한 뒤, phase cut 위치의 표현을 사용하여 다음 페이즈 탈락
위험을 예측한다.

\subsection{예비 예측 실험과 비교 조건}

실험은 맵/플레이 모드 조건을 섞지 않고 따로 나누었다. TPP 스쿼드 5개 맵
(\texttt{Baltic}, \texttt{Tiger}, \texttt{Desert}, \texttt{Neon}, \texttt{Savage})과
FPP 스쿼드 4개 맵(\texttt{Baltic}, \texttt{Tiger}, \texttt{Desert}, \texttt{Neon})으로
총 9개 조건을 만들었다. TPP 조건은 맵마다 300매치, FPP 조건은 맵마다 250매치를
사용했다. 전체 분할은 학습/검증/테스트 \(=2000/250/250\) 매치이고, 각 조건 안에서
동일한 seed 42 분할을 모든 비교 조건에 공유했다. 테스트 요약
기준으로 보류 집합은 약 3.38k개의 페이즈별 생존 팀 라벨을 평가하며, 다음 페이즈
탈락 비율은 약 31.5\%이다.

모든 조건은 같은 다음 페이즈 탈락 위험 예측 헤드를 사용하고, 마스킹된 이진
교차엔트로피(binary cross-entropy, BCE) 손실로 학습하고 평가했다. 모델이 팀
\(i\), 페이즈 \(k\)에 대해 출력한 logit을 \(z_{i,k}\)라고 하면, 다음 페이즈 탈락
확률은 \(\hat{p}_{i,k}=\sigma(z_{i,k})\)이다. \(m_{i,k}\)는 팀 \(i\)가 페이즈 \(k\)
종료 시점까지 생존해 다음 페이즈 라벨이 정의된 경우에만 1이 되는 mask이다. 테스트
손실은 다음과 같이 유효한 팀-페이즈 셀에 대해서만 평균냈다.
\[
\mathcal{L}_{\mathrm{BCE}}
=
-\frac{1}{\sum_{i,k} m_{i,k}}
\sum_{i,k} m_{i,k}
\left[
y_{i,k}\log \hat{p}_{i,k}
+(1-y_{i,k})\log (1-\hat{p}_{i,k})
\right].
\]
따라서 이 값은 정확도의 반대말이 아니라, 실제 탈락 여부에 모델이 부여한 확률을
평가하는 평균 로그 손실이다. 탈락한 팀에 높은 확률을 주고, 생존한 팀에 낮은 탈락
확률을 줄수록 손실이 작아진다.

최신 상태 기준모델은 각 페이즈 종료 시점에서 가장 최근 스냅샷만 인코딩한다. 시간
이력 모델은 팀별 인과적 시간 합성곱을 적용한다. 모든 모델은 1 epoch, batch size 8, Adam
optimizer(lr \(=10^{-3}\), weight decay \(=10^{-4}\)), dropout 0.1, gradient clipping
5.0으로 학습했다. 이 설정은 최종 성능 최적화가 아니라, 시공간 데이터베이스 표현의
각 구성요소가 페이즈별 생존 예측에서 어떻게 작동하는지 확인하기 위한 초기 비교
실험이다.
